% GENERATED FILE. Edit canonical lessons, then run npm run generate:books.
% canonical-source-hash: fnv1a64:0a0124852510d6d5
% canonical-lessons: UR-C50-gana, UR-C50-nachna, UR-C50-milna, UR-C50-rahna, UR-C50-sikhna

\chapter{To sing, To dance, To meet, To live, To learn}
\label{ch:ur-to-sing-to-dance-to-meet-to-live-to-learn}
\hlchaptermodality{50}

\begin{chapteropening}
\textbf{By the end of this chapter:} \emph{I can say to sing, to dance, to meet, to live, to stay and to learn.}

Complete the last lesson of chapter 50: I can say to sing, to dance, to meet, to live, to stay and to learn.
\end{chapteropening}

\section[\texorpdfstring{\ur{گانا} (gānā)}{گانا (gānā)}]{\ur{گانا} (gānā) — to sing; a song}

\label{lesson:UR-C50-gana}



\begin{quote}
Before the new one: say the Urdu for to laugh, then the Urdu for to cry.
\end{quote}

\subsection*{You'll want to know: \ur{گانا}}
\textbf{\ur{گانا}} — \emph{gānā} — \textquotedblleft{}to sing; a song\textquotedblright{}.

\begin{cousinweb}
\textbf{\ur{گانا}} (\emph{gānā}) is a tadbhava, a word inherited from Sanskrit and reshaped by sound change. CDIAL lists it under \emph{gāpayati}, the second form it gives beside \emph{gāyati}, \textquotedblleft{}sings\textquotedblright{} (CDIAL 4135). The noun \textquotedblleft{}a song\textquotedblright{} is the same word, used for what is sung.
\end{cousinweb}

\begin{grammarlens}[title={What you've built}]
One more word for this chapter.
\end{grammarlens}

\subsection*{Your turn}
\begin{itemize}
\raggedright
  \item \emph{Say it:} \emph{gānā}
  \item \emph{Say it:} \emph{gānā}, once more
  \item \emph{Say it:} say \emph{gānā}
  \item \emph{Recall:} say the Urdu for to laugh, then the Urdu for to cry, then say \emph{gānā} again
\end{itemize}

\begin{tcolorbox}[breakable,colback=teal!4,colframe=teal!35!black,title={Before you move on}]
What does \ur{گانا} mean? (\textquotedblleft{}To sing; a song\textquotedblright{}.) Say it once more.
\end{tcolorbox}
\section[\texorpdfstring{\ur{ناچنا} (nāchnā)}{ناچنا (nāchnā)}]{\ur{ناچنا} (nāchnā) — to dance}

\label{lesson:UR-C50-nachna}



\begin{quote}
Before the new one: say the Urdu for to cry, then the Urdu for to sing; a song.
\end{quote}

\subsection*{You'll want to know: \ur{ناچنا}}
\textbf{\ur{ناچنا}} — \emph{nāchnā} — \textquotedblleft{}to dance\textquotedblright{}.

\begin{cousinweb}
\textbf{\ur{ناچنا}} (\emph{nāchnā}) is a tadbhava, a word inherited from Sanskrit and reshaped by sound change. It comes from \emph{nṛtyati}, \textquotedblleft{}dances\textquotedblright{} (CDIAL 7583).
\end{cousinweb}

\begin{grammarlens}[title={What you've built}]
One more word for this chapter.
\end{grammarlens}

\subsection*{Your turn}
\begin{itemize}
\raggedright
  \item \emph{Say it:} \emph{nāchnā}
  \item \emph{Say it:} \emph{nāchnā}, once more
  \item \emph{Say it:} say \emph{nāchnā}
  \item \emph{Recall:} say the Urdu for to cry, then the Urdu for to sing; a song, then say \emph{nāchnā} again
\end{itemize}

\begin{tcolorbox}[breakable,colback=teal!4,colframe=teal!35!black,title={Before you move on}]
What does \ur{ناچنا} mean? (\textquotedblleft{}To dance\textquotedblright{}.) Say it once more.
\end{tcolorbox}
\section[\texorpdfstring{\ur{ملنا} (milnā)}{ملنا (milnā)}]{\ur{ملنا} (milnā) — to meet}

\label{lesson:UR-C50-milna}



\begin{quote}
Before the new one: say the Urdu for to sing; a song, then the Urdu for to dance.
\end{quote}

\subsection*{You'll want to know: \ur{ملنا}}
\textbf{\ur{ملنا}} — \emph{milnā} — \textquotedblleft{}to meet\textquotedblright{}.

\begin{cousinweb}
\textbf{\ur{ملنا}} (\emph{milnā}) is a tadbhava, a word inherited from Sanskrit and reshaped by sound change. It comes from \emph{milati}, \textquotedblleft{}meets\textquotedblright{} (CDIAL 10133).
\end{cousinweb}

\begin{grammarlens}[title={What you've built}]
One more word for this chapter.
\end{grammarlens}

\subsection*{Your turn}
\begin{itemize}
\raggedright
  \item \emph{Say it:} \emph{milnā}
  \item \emph{Say it:} \emph{milnā}, once more
  \item \emph{Say it:} say \emph{milnā}
  \item \emph{Recall:} say the Urdu for to sing; a song, then the Urdu for to dance, then say \emph{milnā} again
\end{itemize}

\begin{tcolorbox}[breakable,colback=teal!4,colframe=teal!35!black,title={Before you move on}]
What does \ur{ملنا} mean? (\textquotedblleft{}To meet\textquotedblright{}.) Say it once more.
\end{tcolorbox}
\section[\texorpdfstring{\ur{رہنا} (rahnā)}{رہنا (rahnā)}]{\ur{رہنا} (rahnā) — to live, to stay}

\label{lesson:UR-C50-rahna}



\begin{quote}
Before the new one: say the Urdu for to dance, then the Urdu for to meet.
\end{quote}

\subsection*{You'll want to know: \ur{رہنا}}
\textbf{\ur{رہنا}} — \emph{rahnā} — \textquotedblleft{}to live, to stay\textquotedblright{}.

\begin{cousinweb}
\textbf{\ur{رہنا}} (\emph{rahnā}) is a tadbhava, a word inherited and reshaped by sound change. It comes from \emph{*rahati}, \textquotedblleft{}remains\textquotedblright{}, a form CDIAL reconstructs from its descendants (CDIAL 10666).
\end{cousinweb}

\begin{grammarlens}[title={What you've built}]
One more word for this chapter.
\end{grammarlens}

\subsection*{Your turn}
\begin{itemize}
\raggedright
  \item \emph{Say it:} \emph{rahnā}
  \item \emph{Say it:} \emph{rahnā}, once more
  \item \emph{Say it:} say \emph{rahnā}
  \item \emph{Recall:} say the Urdu for to dance, then the Urdu for to meet, then say \emph{rahnā} again
\end{itemize}

\begin{tcolorbox}[breakable,colback=teal!4,colframe=teal!35!black,title={Before you move on}]
What does \ur{رہنا} mean? (\textquotedblleft{}To live, to stay\textquotedblright{}.) Say it once more.
\end{tcolorbox}
\section[\texorpdfstring{\ur{سیکھنا} (sīkhnā)}{سیکھنا (sīkhnā)}]{\ur{سیکھنا} (sīkhnā) — to learn}

\label{lesson:UR-C50-sikhna}



\begin{quote}
Before the new one: say the Urdu for to meet, then the Urdu for to live.
\end{quote}

\subsection*{You'll want to know: \ur{سیکھنا}}
\textbf{\ur{سیکھنا}} — \emph{sīkhnā} — \textquotedblleft{}to learn\textquotedblright{}.

\begin{cousinweb}
\textbf{\ur{سیکھنا}} (\emph{sīkhnā}) is a tadbhava, a word inherited from Sanskrit and reshaped by sound change. It comes from \emph{śikṣate}, \textquotedblleft{}learns\textquotedblright{} (CDIAL 12430).
\end{cousinweb}

\begin{grammarlens}[title={What you've built}]
One more word for this chapter.
\end{grammarlens}

\subsection*{Your turn}
\begin{itemize}
\raggedright
  \item \emph{Say it:} \emph{sīkhnā}
  \item \emph{Say it:} \emph{sīkhnā}, once more
  \item \emph{Say it:} say \emph{sīkhnā}
  \item \emph{Recall:} say the Urdu for to meet, then the Urdu for to live, then say \emph{sīkhnā} again
\end{itemize}

\begin{tcolorbox}[breakable,colback=teal!4,colframe=teal!35!black,title={Before you move on}]
What does \ur{سیکھنا} mean? (\textquotedblleft{}To learn\textquotedblright{}.) Say it once more.
\end{tcolorbox}
